\scanpage{33}
\mspage{20}

et l’hyp. de récurrence implique $u_{I'}\in TC$, donc $\widetilde u_{I',I}\in TC$ par coch. de base. D’autre part la flèche $u_I:X_I\mathbin{\unsure{-|\to}}Y_I$ s’identifie à $X_{i_0}\mathbin{-}Y_{i_0}$ ($i_0$ étant élément final de $I$), elle est donc dans $W$. Donc par saturation, $u_{I',I}\in W$, donc $u_{I',I}\in TC$, \unsure{qed.|cqfd.}

Ainsi, la prop.~2 est prouvée, ainsi que son cor.~1 --- et on voit que si dans la prop 2 et son cor., on suppose $I$ \uline{fini}, alors il est inutile dans les hyp pour $C$, resp. $TC=C\cap W$, de supposer la stabilité par composition \uline{transfinie}, la composition ordinaire suffit. Mais il reste \unsure{à} prouver le corollaire 2, que je viens de prouver dans le cas \deleted{\unsure{de}} $I$ \uline{fini}, sans avoir à faire appel à une hypothèse ``transfinie''. Il faut maintenant prouver que c’est \unsure{aussi|encore} OK \deleted{\unsure{si}} dans le cas loc fini, en supposant $C$ stable par \unsure{composition|compositions} \unsure{transfinie|transfinies} \unsure{!} (Il n’y a plus à revenir sur le cas de $W$, \deleted{\unsure{on}} on sait déjà \inserted{maintenant} que $C_I\cap W_{I^\circ}=(TC)_I$, par le cor.~1 déjà prouvé.)
% archive sheet number: 32

